\chapter{Data Manipulation}

\section{\texttt{generate} and \texttt{mutate}}

Create or replace columns. The two forms are equivalent in result; they
differ in mutation semantics:

\begin{center}
\begin{tabular}{lll}
\toprule
\textbf{Form} & \textbf{Semantics} & \textbf{Style} \\
\midrule
\texttt{generate df col = expr}              & Imperative (mutates \texttt{df}) & Stata \\
\texttt{let df2 = generate(df, col = expr)}  & Functional (df unchanged)        & Functional \\
\texttt{df |> generate(col = expr)}          & Pipe assign-back (mutates \texttt{df}) & Chaining \\
\bottomrule
\end{tabular}
\end{center}

\begin{lstlisting}[caption={generate: imperative form}]
load "data.csv" as df

generate df lincome = log(income)
generate df income2 = income^2
generate df dummy   = if(income > 1000, 1, 0)
\end{lstlisting}

\begin{lstlisting}[caption={mutate: functional form}]
let df2 = mutate(df, lincome = log(income), income2 = income^2)
// df does not have lincome or income2
// df2 has them
\end{lstlisting}

\section{\texttt{filter}}

Selects rows that satisfy a condition:

\begin{lstlisting}
let recent = filter(df, year >= 2010)
let high   = filter(df, income > 5000 and state == "TX")
\end{lstlisting}

The imperative form uses \lstinline{keep}:

\begin{lstlisting}
keep df if year >= 2010
\end{lstlisting}

\section{\texttt{select} and \texttt{drop}}

\begin{lstlisting}
// keep only these columns
let slim   = select(df, year, state, income)
let no_id  = drop(df, id, code)    // remove these columns
\end{lstlisting}

\section{\texttt{rename}}

\begin{lstlisting}
rename df old_name new_name
// rename family_income to income
rename df family_income income
\end{lstlisting}

\section{\texttt{sort}}

\begin{lstlisting}
sort df year              // ascending
sort df income desc       // descending
sort df state year desc   // multiple columns
\end{lstlisting}

\section{\texttt{dropna}}

Removes rows with missing values:

\begin{lstlisting}
dropna df                 // remove rows with any NaN
dropna df income gini     // remove rows with NaN in income or gini
\end{lstlisting}

\section{\texttt{append}}

Stacks two DataFrames vertically (equivalent to \texttt{rbind}):

\begin{lstlisting}
let df_total = append(df_2024, df_2025)
\end{lstlisting}

Columns must be compatible. Columns missing from one DataFrame are filled
with \texttt{NaN}.

\section{\texttt{collapse} and \texttt{group\_by}}

\lstinline{group_by} aggregates by groups:

\begin{lstlisting}[caption={Aggregation by group}]
let state_means = group_by(df, state,
  mean_income = mean(income),
  n           = count()
)
\end{lstlisting}

\lstinline{collapse} is the Stata-style imperative form:

\begin{lstlisting}
collapse df (mean) income gini, by(state year)
\end{lstlisting}

\section{\texttt{pivot\_longer} and \texttt{pivot\_wider}}

\begin{lstlisting}[caption={Wide to long}]
let long = pivot_longer(df,
  cols   = ["q1", "q2", "q3"],
  names  = "quarter",
  values = "value"
)
\end{lstlisting}

\begin{lstlisting}[caption={Long to wide}]
let wide = pivot_wider(df,
  id     = "county",
  names  = "year",
  values = "gdp"
)
\end{lstlisting}

\section{\texttt{replace}}

Replaces values conditionally:

\begin{lstlisting}
// zero out negative values
replace df income = 0 if income < 0
replace df gini = . if gini > 1    // mark as missing
\end{lstlisting}

\section{\texttt{tabulate}}

Frequency table:

\begin{lstlisting}
tabulate df state
tabulate df state year    // cross-tabulation
\end{lstlisting}
